\chapter{Einführung}

\section{Übersicht über die Applikation}

\textbf{PICasso} ist ein Simulator für PIC-Midrange-Mikrocontroller, geschrieben in C++17.
Die Applikation liest assemblierte PIC-Programme im \code{.LST}-Format ein und führt diese Instruktion für Instruktion aus.
Dabei werden alle wesentlichen Hardwarekomponenten des PIC simuliert:
ein 2-Bank-RAM mit Registerspiegelung, ein 8-Bit-W-Akkumulator, ein 8-Level-Hardware-Stack, ein Timer/Prescaler-Subsystem
sowie die vollständige Ausführung aller 35 PIC-Midrange-Instruktionen.
Die Applikation bietet eine interaktive Terminal-Oberfläche (ncurses), über die der Benutzer Programme laden,
schrittweise oder im Dauerlauf ausführen, RAM-Inhalte inspizieren und modifizieren sowie einzelne Bits umschalten kann
Die Simulation läuft in einem separaten Thread, während die UI im Hauptthread rendert -- synchronisiert über atomare Flags und Mutexe.

\textbf{Zweck:} PICasso löst das Problem, PIC-Assemblerprogramme ohne physische Hardware testen und debuggen zu können.
Es dient als Lern- und Entwicklungswerkzeug für eingebettete Systeme.

\section{Starten der Applikation}

\textbf{Voraussetzungen:}

\begin{itemize}
    \item C++17-kompatibler Compiler (z.B. g++ $\geq$ 9)
    \item CMake $\geq$ 3.14
    \item ncurses-Bibliothek (\code{libncurses-dev} unter Linux)
\end{itemize}

\textbf{Schritt-für-Schritt-Anleitung:}

\begin{lstlisting}[style=codeblock,language=bash]
# 1. Repository klonen / Projektverzeichnis wechseln
git clone https://github.com/DefinitlyNotABot/PICasso.git
cd PICasso

# 2. Build-Verzeichnis erstellen und CMake konfigurieren
mkdir -p build && cd build
cmake ..

# 3. Projekt kompilieren
make

# 4. Simulator starten
./PICasso

# 5. Tests ausführen (optional)
cd ..
./build/AllTests
\end{lstlisting}

Nach dem Start öffnet sich die ncurses-Oberfläche.
Über den \emph{Load}-Button kann eine \code{.LST}-Datei aus dem \code{progs/}-Verzeichnis geladen werden.
Anschließend kann das Programm mit \emph{Step}, \emph{Run} (Dauerlauf) oder \emph{Dash} (Schnelllauf) ausgeführt werden.

\section{Technischer Überblick}

\begin{longtable}{|p{0.18\textwidth}|p{0.22\textwidth}|p{0.5\textwidth}|}
\hline
\textbf{Technologie} & \textbf{Beschreibung} & \textbf{Begründung} \\
\hline
C++17 & Primäre Programmiersprache & Bietet hardwarenahe Kontrolle, geeignet für Simulator-Entwicklung;
unterstützt \code{std::optional}, \code{std::shared_ptr}, \code{std::filesystem} und strukturierte Bindings. \\
\hline
CMake & Build-System & Plattformunabhängiges Build-System;
ermöglicht automatisches Herunterladen von Abhängigkeiten (Catch2) via \code{FetchContent}. \\
\hline
ncurses & Terminal-UI-Bibliothek & Ermöglicht eine interaktive TUI ohne grafische Abhängigkeiten;
ideal für ein konsolenbasiertes Debugging-Tool. Unterstützt Farben, Mauseingabe und Fenstermanagement. \\
\hline
Catch2 & Unit-Test-Framework & Header-only-Framework mit BDD-Stil (\code{TEST_CASE}/\code{SECTION});
automatische Testregistrierung und Generatoren (\code{GENERATE}). \\
\hline
nlohmann/json & JSON-Parsing für Testdateien & Erlaubt deklarative Testspezifikation in JSON-Dateien;
einfache Integration als Header-only-Bibliothek. \\
\hline
\code{std::thread / std::mutex / std::atomic} & Multithreading & Simulation und UI laufen parallel;
atomare Flags steuern den Simulationszustand. \\
\hline
\end{longtable}

